\documentclass[mscthesis]{usiinfthesis}
\usepackage{lipsum}
\usepackage{float}

\usepackage{listings}
\lstset{
    language=Python,
    basicstyle=\ttfamily,
    keywordstyle=\bfseries,
    showstringspaces=false
}

\lstdefinelanguage{algebra}
{morekeywords={import,sort,constructors,observers,transformers,axioms,if,
else,end},
sensitive=false,
morecomment=[l]{//s},
}


\setlength{\parindent}{0pt}

\title{Benchmark Unlearning for Fair Code LLM Evaluation}
\author{Raffaele Perri}
\begin{committee}
\advisor{Prof.}{Gabriele}{Bavota}
\coadvisor{}{Alessandro}{Giagnorio}{}
\end{committee}
\Day{Yesterday} %compulsory
\Month{June} %compulsory
\Year{2026}
\place{Lugano} %compulsory

\dedication{To my beloved} %optional
\openepigraph{Someone said \dots}{Someone} %optional

% \makeindex %optional, also comment out \theindex at the end
\begin{document}

\maketitle %generates the titlepage, this is FIXED

\frontmatter %generates the frontmatter, this is FIXED

\begin{abstract}
This is a very abstract abstract. 

\lipsum
\end{abstract}

\tableofcontents 
\listoffigures %optional
\listoftables %optional

\mainmatter

\chapter{Introduction}
 
\chapter{Background and Related Work}
This chapter establishes the theoretical and empirical foundations of the thesis. \\
The first section describes how large language models are evaluated on coding tasks and what constitutes a benchmark in this context. The second introduces the core concepts of machine unlearning. The final section, Related Work, surveys existing code generation benchmarks, unlearning approaches, and empirical studies on data leakage and benchmark contamination in code LLMs — together motivating the targeted, unlearning-based decontamination this thesis investigates.

\section{Code Generation Benchmarks}
A benchmark is a standardized collection of tasks — typically pairs of problem specifications and reference solutions or automated test suites — that allows the research community to compare the coding competence of different models under controlled, reproducible conditions. In code generation, a benchmark generally consists of a prompt (a natural-language description or partial code stub), an expected output (a function or program), and one or more test cases acting as an oracle for functional correctness. The evaluation score is computed by running model-generated code against these test cases.\\

The standard evaluation pipeline proceeds as follows. The model receives a problem prompt — which may include a function signature, a docstring, and optional usage examples:
\begin{lstlisting}
	def is_palindrome(s: str) -> bool:
		""" Return True if the given string s is a palindrome,
		i.e., it reads the same forwards and backwards.
		Ignore case differences.
		>>> is_palindrome("racecar")
		True
		"""
\end{lstlisting}
It is asked to produce a complete implementation:
\begin{lstlisting}
	s = s.lower()
	return s == s[::-1]
\end{lstlisting}

The generated code is then executed in a sandboxed environment against a predefined set of test cases. A solution is marked correct if all tests pass:
\begin{lstlisting}
	assert is_palindrome("racecar") == True
	assert is_palindrome("hello") == False
	assert is_palindrome("Madam") == True
	assert is_palindrome("") == True
	assert is_palindrome("A") == True
\end{lstlisting}

The primary metric for measuring functional correctness in this setting is $pass@k$. \\
For each problem in the benchmark, the model generates $k$ independent code samples; the problem is considered solved if at least one sample passes all test cases. The final score is the fraction of benchmark problems satisfying this condition. 
Intuitively, $pass@1$ measures whether the model produces a correct solution on the first attempt, while $pass@10$ or $pass@100$ measure whether a correct solution appears within a larger pool of attempts — a useful signal when evaluating more diverse or stochastic generation strategies.\\

Although the $pass@k$ metric predates HumanEval, \citet{DBLP:journals/corr/abs-2107-03374} identified a flaw in its naive formulation: generating exactly $k$ samples and checking whether any pass produces high-variance estimates. They proposed an unbiased estimator — generate $n \geq k$ samples per problem, count the number $c$ that pass all tests, and compute $pass@k$ analytically via a closed-form expression. This estimator has since become standard across code generation benchmarks. In practice, $pass@1$ computed with greedy or low-temperature sampling is the most commonly reported variant, as it directly reflects the probability that a single generation attempt produces a correct solution.\\

This execution-based evaluation pipeline is the standard approach across code generation benchmarks. It provides an objective, binary signal of correctness — either the code passes all test cases or it does not — avoiding the ambiguity of surface-level textual metrics such as BLEU (\citet{BLUEamethodforautomaticevaluationofmachinetranslation}) or CodeBLEU (\citet{ren2020codebleumethodautomaticevaluation}). These metrics were originally designed for natural language and may assign high scores to syntactically similar but functionally incorrect programs. \citet{Evtikhiev_2023} confirm this limitation empirically, showing that BLEU and CodeBLEU correlate poorly with human judgments of code correctness and can rank functionally incorrect outputs above correct ones.\\

Benchmark evaluation also requires careful test suite design. Inadequate or trivially simple tests can yield inflated accuracy figures that do not reflect true model capability (\citet{zhang2024naturalcodebenchexaminingcodingperformance}; \citet{ouyang2025dscodebenchrealisticbenchmarkdata}; \citet{koohestani2025benchmarkingaimodelssoftware}). For instance, \citet{ouyang2025dscodebenchrealisticbenchmarkdata} reports that DS-1000 provides an average of only 2.1 test cases per problem — a limitation that motivates DSCodeBench's design of 200 test cases per problem to ensure broader coverage of code behavior, including edge cases and intermediate states.\\

A further threat to benchmark validity is data contamination: benchmark problems or their solutions appear in a model's pretraining or fine-tuning corpus, causing evaluation scores to reflect memorization rather than genuine generalization. This is a particularly acute problem in the code domain, as widely used benchmarks such as HumanEval and MBPP are publicly available across numerous repositories and may have been seen by a model prior to evaluation. Contaminated models perform significantly better on problems encountered during training than on genuinely unseen ones (\citet{riddell2024quantifyingcontaminationevaluatingcode}; \citet{matton2024leakagecodegenerationevaluation}; \citet{zhou2025lessleakbenchinvestigationdataleakage}). These findings motivate the central goal of this thesis: rather than preventing contamination at the data ingestion stage, this work aims to neutralize it post-hoc by identifying and suppressing the model neurons responsible for retaining benchmark-specific knowledge.

\section{Machine Unlearning}
Understanding machine unlearning requires first understanding how LLMs acquire knowledge. Modern LLMs are trained in two main phases (\citet{liu2024rethinkingmachineunlearninglarge}).  In the first phase, called pre-training, the model is exposed to a huge amount of raw text and code; trillions of tokens drawn from online resources, repositories, books. Then, it is trained to predict the next token in a sequence. Through this process, the model encodes: world knowledge, language patterns, and coding conventions into its parameters. Any information present in the pre-training corpus may be memorized, including benchmark problems and their solutions. 
\par\vspace{1em}
In the second phase, the pre-trained model is fine-tuned on smaller, curated datasets to specialize it for a particular task or to align it with human preferences. This takes the form of supervised fine-tuning (SFT) on instruction-response pairs, or reinforcement learning from human feedback (RLHF), which uses human ratings to steer outputs toward being more helpful, harmless, and honest (\citet{ouyang2022traininglanguagemodelsfollow}). For code-specific models, fine-tuning typically involves additional training on curated code datasets, commit histories, or synthesized instruction-following examples. This pipeline is now standard for building capable code LLMs such as CodeLlama, DeepSeek-Coder, and Qwen-Coder (\citet{rozière2024codellamaopenfoundation}; \citet{ouyang2025dscodebenchrealisticbenchmarkdata}). 
\par\vspace{1em}
The knowledges absorbed during the pre-training phase are embedded in the weights of the model and it is not easy to remove them.
Fine-tuning on new data is not a  guarantee that previously memorized information is erased — it may be partly suppressed , remaining recoverable under certain prompting conditions (\citet{hong2024dissectingfinetuningunlearninglarge}). This is the problem that machine unlearning addresses.
\par\vspace{1em}
Machine unlearning (MU) refers to the specific removal of some training data or knowledge from a trained model, that were influencing the model, without retraining from scratch. Originally motivated by data privacy regulations (GDPR's right to be forgotten), MU has expanded to address copyright removal, safety alignment, bias mitigation, and hallucination reduction in LLMs (\citet{liu2024rethinkingmachineunlearninglarge}).
\par\vspace{1em}
Exact unlearning — retraining from scratch on the dataset with the target data removed — is the most reliable approach but is computationally infeasible for large models. Research has therefore concentrated on approximate unlearning methods that modify a trained model post-hoc. These methods must balance two competing objectives: unlearning efficacy, meaning that targeted knowledge is effectively erased, and utility preservation, meaning that general model capabilities outside the unlearning scope remain intact.
\par\vspace{1em}
Gradient Ascent (GA) is one of the most know unlearning methods: updates parameters of the model by maximizing the likelihood of mis-predictions for the samples in the forget set (\citet{yao2024largelanguagemodelunlearning}). However, during the optimization, GA alone can be sensitive to the choice of hyperparameters, which can lead to unlearning failures as catastrophic collapse  (\citet{liu2024rethinkingmachineunlearninglarge}). This has given rise to several improved variants.
\par\vspace{1em}
Gradient Difference (GradDiff) is the combination of GA on the forget set with gradient descent on the retain set and a KL divergence regularization term to minimize the divergence between the unlearned model and the original model on a reference dataset, preventing catastrophic deterioration of general capability (\citet{hong2024dissectingfinetuningunlearninglarge}; \citet{wang2024llmunlearninglossadjustment}).
\par\vspace{1em}
Negative Preference Optimization (NPO) treats the forgotten data exclusively as negative examples: it turns the unlearning problem into a minimization problem mitigating the catastrophic collapse compared to GA (\citet{liu2024rethinkingmachineunlearninglarge}; \citet{jia2025waglestrategicweightattribution}).
\par\vspace{1em}
Preference Optimization (PO)  takes the opposite direction from NPO: instead of treating forget data as negative examples, it introduces targeted unlearning responses ("I don't know") as positive examples, producing a bounded forget loss (\citet{jia2025waglestrategicweightattribution}).
Another variant transforms GA into a gradient descent approach by minimizing the likelihood of predictions on relabeled forgetting data, replacing the original responses with generic or uninformative alternatives (\citet{yao2024largelanguagemodelunlearning}; \citet{liu2024rethinkingmachineunlearninglarge}).
\par\vspace{1em}
A related approach is Task Arithmetic (\citet{ilharco2023editingmodelstaskarithmetic}; \citet{cai2026perparametertaskarithmeticunlearning}): a task vector is defined as the element-wise difference between a fine-tuned model's weights and the original pre-trained model's weights.
Negating this vector and subtracting it from the pre-trained model degrades performance on the target task while leaving unrelated capabilities intact.
\par\vspace{1em}
A distinct line of work frames unlearning through mechanistic interpretability. Rather than modifying all model parameters globally, these methods first localize the neurons or weight subsets responsible for storing the target knowledge, then restrict modifications to those components. This limits collateral damage to unrelated capabilities and improves both unlearning efficacy and utility preservation. This neuron-based paradigm is the foundational principle of the approach developed in this thesis.

\section{Related Work}
Below, we present a survey of previous studies related to code generation benchmarks, machine unlearning techniques, and data leakage in code LLMs.

\subsection{Code Generation Benchmarks}
HumanEval (\citet{DBLP:journals/corr/abs-2107-03374}) is arguably the most widely adopted benchmark for evaluating code generation. It consists of 164 hand-crafted Python problems, each defined by a function signature, a docstring, and a hidden test suite. Solutions are evaluated via execution against unit tests. Its companion, MBPP (Mostly Basic Python Problems) (\citet{austin2021programsynthesislargelanguage}), is a crowdsourced collection of 974 Python programming problems of varying difficulty with a natural-language description and on average three test cases (\citet{ouyang2025dscodebenchrealisticbenchmarkdata}; \citet{koohestani2025benchmarkingaimodelssoftware}). Despite their popularity — virtually all model announcements reporting code capability cite at least one of them — both benchmarks suffer from well-documented limitations: they are narrowly Python-centric, contain predominantly algorithmic and basic programming problems (96.9\% and 89.5\% of their content, respectively  (\citet{zhang2024naturalcodebenchexaminingcodingperformance})), and provide limited test coverage that makes evaluation scores statistically unstable (\citet{koohestani2025benchmarkingaimodelssoftware}).
To overcome the limitations of HumanEval and MBPP, a range of more comprehensive and challenging benchmarks has proposed by the research community.
\par\vspace{1em}
NaturalCodeBench (NCB) (\citet{zhang2024naturalcodebenchexaminingcodingperformance}) contains 402 problems drawn from real coding online services across six domains (software engineering, data science, algorithms, systems administration, AI systems, and front-end development).
Differs from HumanEval and MBPP since it features multi-file format and it is evaluated in an executable Docker environment, to have a more realistic setting.
Notably, a performance mismatch study showed that many models ranked highly on HumanEval perform substantially worse on NCB, suggesting overfitting to the former.
\par\vspace{1em}
DSCodeBench (\citet{ouyang2025dscodebenchrealisticbenchmarkdata}) contains 1,000 problems sourced from real GitHub repositories and covering ten popular Python libraries (NumPy, Pandas, SciPy, Scikit-learn, TensorFlow, PyTorch, Matplotlib, Seaborn, Keras, and LightGBM). and targets the data science domain specifically. Each problem has a solution code of 22.5 lines and a rich test suites of 200 cases, on average, providing significantly stronger evaluation than prior data science benchmarks such as DS-1000.
\par\vspace{1em}
Other notable specialized benchmarks include: BigCodeBench (\citet{zhuo2025bigcodebenchbenchmarkingcodegeneration}), that targets diverse function calls across 62 libraries; SWE-Bench (\citet{jimenez2024swebenchlanguagemodelsresolve}), that evaluates practical software engineering by asking models to fix real GitHub issues; EditBench (\citet{chi2025editbenchevaluatingllmabilities}), that assesses instructed code editing over 540 real-world problems; GitChameleon 2.0 (\citet{misra2025gitchameleon20evaluatingai}), that evaluates version-conditioned code generation across Python library updates; and LiveCodeBench (\citet{jain2024livecodebenchholisticcontaminationfree}), a dynamic benchmark that continuously aggregates new problems from competitive programming platforms specifically to avoid overfitting and reduce the risk of data contamination (\citet{koohestani2025benchmarkingaimodelssoftware}). For multi-language evaluation, benchmarks such as HumanEval-X (\citet{zheng2024codegeexpretrainedmodelcode}; \citet{koohestani2025benchmarkingaimodelssoftware}), MultiPL-E (\citet{cassano2022multiplescalableextensibleapproach}; \citet{koohestani2025benchmarkingaimodelssoftware}), XLCoST (\citet{koohestani2025benchmarkingaimodelssoftware}), and McEval (\citet{chai2024mcevalmassivelymultilingualcode}) do not contain only Python-centric problems but extend their scope to dozens of programming languages. McEval is of particular relevance to this thesis, as our work builds upon an improved version of it.
\par\vspace{1em}
A comprehensive survey (\citet{koohestani2025benchmarkingaimodelssoftware}) catalogues hundreds of code benchmarks organized along multiple dimensions including: task type (generation, summarization, translation, repair, editing), programming language, and evaluation protocol. This reveals that while the field has made significant improvements  in benchmark diversity, it still suffers from recurring issues: shallow test coverage, limited language support, and proximity to training data, undermining the reliability of evaluations.

\subsection{Machine Unlearning}
An increasing number of work has explored machine unlearning for LLMs. \citet{yao2024largelanguagemodelunlearning} introduce a gradient ascent-based framework targeting harmful and privacy-violating content in LLMs. \citet{liu2024rethinkingmachineunlearninglarge} provide a comprehensive rethinking of the field, surveying and unifying gradient-based, PEFT-based, localization-informed, and influence-function approaches under a common framework. \citet{wang2024llmunlearninglossadjustment} propose FLAT, a loss adjustment method that performs unlearning using only the forget data, without requiring access to a retain set or reference model. \citet{hong2024dissectingfinetuningunlearninglarge} dissect the internal mechanics of fine-tuning-based unlearning, showing through activation patching and parameter restoration experiments that these methods alter how knowledge is retrieved rather than truly erasing it. \citet{ilharco2023editingmodelstaskarithmetic} demonstrate that negating a fine-tuning task vector can effectively suppress a learned capability while largely preserving general model performance.
\par\vspace{1em}
While these approaches cover a broad spectrum of unlearning techniques, the focus of this thesis is on neuron-based and localization-informed methods, which offer the most targeted and precise form of knowledge removal. Rather than applying gradient updates across all model parameters — which risks degrading unrelated capabilities — these approaches first identify the specific model components responsible for storing the target knowledge, and then restrict modifications to those components only. We now review the most relevant works in this line.
\par\vspace{1em}
\citet{kojima2024multilingualabilitydecoderbasedpretrained} provide the foundational insight for this thesis. They demonstrate that language-specific behaviors in multilingual LLMs can be attributed to a sparse subset of neurons in the feed-forward network layers. By analyzing activation patterns conditioned on language-specific inputs, they identify these neurons and show that patching (zeroing out) them at inference time achieves targeted capability removal without broadly degrading model performance. 
\par\vspace{1em}
\citet{pochinkov2024dissectinglanguagemodelsmachine} introduce a post-hoc approach where, to identify neurons contributing to a specific capability (e.g. coding ability or toxic behavior), they used the influence functions computed after training  — and to remove those neurons while leaving the rest of the network intact they used unstructured pruning. This work demonstrates that capability-level unlearning can be achieved through surgical weight removal without any gradient-based optimization.
\par\vspace{1em}
\citet{jia2025waglestrategicweightattribution} propose WAGLE, which extends the localization idea to weight-level attribution using bi-level optimization. Rather than relying on simple magnitude-based pruning or gradient-based saliency, WAGLE derives closed-form attribution scores that jointly consider the unlearning objective and the utility retention objective. Experiments on the TOFU and WMDP benchmarks show that restricting gradient updates to WAGLE-selected weights (typically 80-95\% density) consistently improves unlearning efficacy while maintaining utility, outperforming dense model-based baselines.
\par\vspace{1em}
\citet{cai2026perparametertaskarithmeticunlearning} further refine the task vector negation approach through per-parameter scaling. They address a key limitation of naive task vector subtraction: because the forget target and other knowledge may be coupled in the same parameters, a uniform negation risks corrupting retained capabilities. Their per-parameter scaling allows finer-grained control over which parts of the model are affected by the unlearning operation.
\par\vspace{1em}
Turning to unlearning applied specifically to code language models, \citet{chu2026scrubitouterasingsensitivememorization} address the problem of sensitive memorization — including leaked API keys, email addresses, and IP addresses. They introduce CodeEraser, a selective unlearning framework that performs gradient ascent on sensitive segments of code while applying KL divergence regularization on non-sensitive contexts, minimizing collateral damage to the model's broader coding utility. Results on CodeParrot, CodeGen, and InCoder show that memorized sensitive content can be effectively erased while preserving benchmark performance on HumanEval.
\par\vspace{1em}
\citet{jiang2025largelanguagemodelunlearning} focus on unlearning copyrighted code and highlight the tension between removing memorized content and preserving coding ability. This work underscores that, in the code domain, the granularity and precision of the unlearning target are especially important: overly broad unlearning degrades general-purpose coding capabilities, while overly narrow unlearning may not fully suppress memorized knowledge — a challenge that neuron-level targeting is well-positioned to address.

\subsection{Data Leakage in Code LLMs}
\citet{Sallou_2024} argue that researchers using LLMs in software engineering experiments must contend with three intertwined threats to validity: the opacity of closed-source models' training data, the blurry separation between training and evaluation sets, and the reproducibility of results across model versions. They provide a striking illustration: ChatGPT-3.5, though not fine-tuned for bug-fixing tasks, exhibits precise knowledge of specific bugs in the Defects4J dataset — almost certainly because descriptions of those bugs appeared in its pretraining corpus. This foreknowledge invalidates construct validity (the model is evaluated on data it has seen) and external validity (results may not generalize to unknown code).
\par\vspace{1em}
\citet{riddell2024quantifyingcontaminationevaluatingcode} perform a systematic study of overlap between popular code generation benchmarks and two large pretraining corpora (The Pile and The Stack). Using both surface-level string matching and semantic-level embedding similarity, they find substantial overlap: 12.2\% of HumanEval problems are present in The Pile, and 18.9\% in The Stack; for MBPP the figures are 6.9\% and 32.2\%, respectively. Critically, models perform significantly better on the contaminated subset of problems than on the non-contaminated subset, confirming that contamination inflates reported scores.  By analyzing the factors that modulate effects of contamination, emerged that larger models and shorter problems are more susceptible to memorization.
\par\vspace{1em}
\citet{matton2024leakagecodegenerationevaluation} distinguish three contamination ways: (1) direct leakage, where benchmark problems, both in original or paraphrased forms, are directly present in training corpora (e.g. every HumanEval prompt appears at least 43 times in public GitHub repositories); (2) indirect leakage, where large-scale code synthesis pipelines produce training examples highly similar to benchmark problems,  through synthetic data; and (3) overfitting through checkpoint selection, where repeated evaluation on public benchmarks during development leads to selection bias towards model versions that perform well on leaked data. The authors introduce LBPP (Less Basic Python Problems), a benchmark of 161 carefully crafted problems designed to be free of contamination, and demonstrate that leading models achieve pass@1 scores up to 43% lower on LBPP than on HumanEval — a striking indication of the extent to which reported HumanEval performance reflects overfitting rather than true coding capability.
\par\vspace{1em}
\citet{zhou2025lessleakbenchinvestigationdataleakage} conduct the most comprehensive leakage audit to date, examining data leakage across 83 software engineering benchmarks (spanning: Python, Java, C, C++, and other languages). The methodology combines automatic detection (MinHash + Locality-Sensitive Hashing on n-gram representations of code) with manual verification by experienced developers to eliminate false positives. The study finds an average leakage ratio of 4.8\% across Python SE benchmarks, but with striking variation: QuixBugs exhibits 100\% leakage (all 40 problems appear in pretraining data), while SWE-Bench shows 8.7\% leakage. Notably, even benchmarks considered 'safe' such as HumanEval (1.8\%) and MBPP (0.4\%) show non-negligible overlap.
\par\vspace{1em}
Several methods have been proposed to detect whether a model has been trained on a given benchmark. \citet{shi2024detectingpretrainingdatalarge} propose a Min-K\% Prob method that exploits the observation that a model assigns higher minimum token probabilities to memorized text than to novel text. \citet{zawalski2025detectingdatacontaminationllms} show that a model's ability to complete partially revealed benchmark problems — using only a few in-context examples as cues — serves as evidence of prior exposure. \citet{stoehr2024localizingparagraphmemorizationlanguage} further show that memorization of multi-sentence passages can be traced to specific, localized regions of the model via activation analysis, providing mechanistic evidence for the viability of targeted unlearning.
\par\vspace{1em}
Despite these advances in detection, removal of detected contamination remains an open problem. Filtering training corpora is error-prone and incomplete; retraining is expensive; and alternative benchmark creation is a race that will eventually be lost as models train on ever-larger internet scrapes. These limitations collectively motivate the approach explored in this thesis: rather than preventing contamination upstream, we aim to neutralize it downstream by identifying and zeroing out the model neurons responsible for storing benchmark-specific knowledge, thereby restoring the fairness of benchmark-based evaluation without the cost of retraining.



\appendix %optional, use only if you have an appendix

\backmatter

\chapter{Glossary} %optional

%\bibliographystyle{alpha}
%\bibliographystyle{dcu}
\bibliographystyle{plainnat}
\bibliography{biblio}

%\cleardoublepage
%\theindex %optional, use only if you have an index, must use
	  %\makeindex in the preamble
\lipsum

\end{document}
